\documentclass[12pt]{report}


% --- Шрифты и локализация (XeLaTeX) ---
\usepackage{fontspec}
\usepackage[russian]{babel}

\setmainfont{Liberation Serif}
\setsansfont{Liberation Sans}
\setmonofont{Liberation Mono}

% --- Математические пакеты (загружаем до tcolorbox) ---
\usepackage{amsmath,amssymb,amsfonts}

% --- Настройка полей страницы ---
\usepackage{geometry}
\geometry{
    a4paper,
    top=25mm,
    bottom=25mm,
    left=30mm,
    right=20mm,
    headheight=25pt,   % Устраняет предупреждение fancyhdr
    footskip=10mm
}

% --- Цвета и заголовки ---
\usepackage{xcolor}
\usepackage{titlesec}

\definecolor{mainblue}{RGB}{26, 82, 118}   
\definecolor{accentblue}{RGB}{41, 128, 185} 
\definecolor{graybg}{RGB}{245, 247, 250}   

\titleformat{\chapter}[display]
  {\normalfont\huge\bfseries\color{mainblue}}
  {\filleft\normalfont\LARGE\sffamily\color{accentblue}\chaptertitlename\ \thechapter}
  {20pt}
  {\filleft\Huge}
  [\vspace{1ex}\titlerule]

\titleformat{\section}
  {\normalfont\Large\bfseries\color{mainblue}\sffamily}
  {\thesection}{1em}{}

\titleformat{\subsection}
  {\normalfont\large\bfseries\color{accentblue}\sffamily}
  {\thesubsection}{1em}{}

% --- Блоки tcolorbox ---
\usepackage{tcolorbox}
\tcbuselibrary{theorems}

% Блок для важных заметок
\newtcolorbox{note}[1][Важное примечание]{
    colback=graybg,
    colframe=accentblue,
    fonttitle=\bfseries\sffamily,
    coltitle=white,
    title=#1,
    arc=4mm,
    boxrule=0.5mm,
    before skip=10pt,
    after skip=10pt
}

% Блок для определений
\newtcolorbox{definition}[1][Определение]{
    boxrule=0pt,
    leftrule=3mm,
    rightrule=0pt,
    toprule=0pt,
    bottomrule=0pt,
    colframe=mainblue,
    colback=graybg,
    fonttitle=\bfseries\sffamily\color{mainblue},
    title=#1,
    attach title to upper,
    after title={\par\smallskip},
    arc=0mm,
    before skip=10pt,
    after skip=10pt
}

% Блок для теорем
\newtcbtheorem[number within=chapter]{theorem}{Теорема}{
    boxrule=0pt,
    leftrule=3mm,
    rightrule=0pt,
    toprule=0pt,
    bottomrule=0pt,
    colframe=accentblue,
    colback=graybg,
    fonttitle=\bfseries\sffamily\color{mainblue},
    coltitle=mainblue,
    attach title to upper,
    after title={\par\smallskip},
    arc=0mm,
    before skip=12pt,
    after skip=12pt
}{th}

% --- Межстрочный интервал и абзацы ---
\usepackage{setspace}
\onehalfspacing
\setlength{\parindent}{1.25cm}
\setlength{\parskip}{4pt}

% --- Списки ---
\usepackage{enumitem}
\setlist[itemize]{label=\color{accentblue}$\bullet$, topsep=4pt, itemsep=2pt}
\setlist[enumerate]{label=\color{mainblue}\arabic*., topsep=4pt, itemsep=2pt}

% --- Колонтитулы ---
\usepackage{fancyhdr}
\pagestyle{fancy}
\fancyhf{}
\fancyhead[L]{\small\sffamily\color{gray}\leftmark}
\fancyfoot[C]{\small\sffamily\thepage}
\renewcommand{\headrulewidth}{0.4pt}
\renewcommand{\headrule}{\hbox to\headwidth{\color{accentblue}\leaders\hrule height \headrulewidth\hfill}}

% --- Графика и гиперссылки ---
\usepackage{graphicx}
\usepackage{hyperref}
\usepackage{bookmark} % Ускоряет генерацию outline и убирает rerunfilecheck warnings
\hypersetup{
    colorlinks=true,
    linkcolor=mainblue,
    filecolor=magenta,      
    urlcolor=accentblue,
    pdftitle={Учебное пособие},
    pdfencoding=auto
}

% --- Метаданные ---
\author{Д. М. Востриков}
\title{Формальная верификация в математике}
\date{\the\year}

% =======================================================
\begin{document}
\maketitle
\tableofcontents

\newpage

\chapter{Глава 1}
\section{Лекция 1. Бестиповое \texorpdfstring{$\lambda$}{λ}-исчисление}

\begin{definition}
Бестиповое $\lambda$-исчисление --- это формальная система для изучения функций и вычислений.
\end{definition}

Пусть $x \in \mathbb{R}$:
\begin{itemize}
    \item $\sin{1}$ --- действительное число;
    \item $\sin{x}$ --- также действительное число.
\end{itemize}

Форма записи $(\lambda x.(\sin{x}+\cos{x}))(\pi)$: взяли функцию и применили её к числу (аппликация) $\sin{\pi} + \cos{\pi}$.

Фиксируем множество:
\begin{itemize}
    \item $Var = \{x, y, z, \dots\}$
    \item $\Lambda$ --- множество $\lambda$-термов
    \item $Var \subset \Lambda$
    \item $M, N \in \Lambda \implies (M \, N \in \Lambda)$
    \item $x \in Var, \; M \in \Lambda \implies (\lambda x. M) \in \Lambda$
\end{itemize}

Соглашения: 
\[ x \, y \, z := (xy)z \]
\[ \lambda x. \lambda y. M = \lambda xy. M \]

\begin{theorem}{Теорема Черча — Россера}{cr1}
$\exists K$ такое, что если $N$ и $L$ $\beta$-эквивалентны $K$, то $N$ и $L$ также являются $\beta$-редукциями $M$.
\end{theorem}

\begin{theorem}{Теорема Черча — Россера (следствие)}{cr2}
Если у терма $M$ существует нормальная форма $N$, то к ней всегда можно прийти, используя нормальную стратегию редукции.
\end{theorem}

\[ True := \lambda x y. x \]
\[ False := \lambda x y. y \]
\[ And := (\lambda x y. x y)\, False \]
\[ Not := \lambda x. x \, False \, True \]

\end{document}